%%%PAGE 300 rot=0 legibility=good summary=动态圆（放缩圆）示意图：过同一点的一族圆，红笔标注弦最长、直径当弦长、圆心角与弦切角的关系及最长时间公式，下半页空白
%%%BLOCK id=300-01 ch=11 sec=有界磁场·动态圆 kind=diagram alt=none cont=prev y=0.00-0.50
% 承接上一页(p299)“弦切角”处延伸过来的红色长曲线：自本页左缘 y≈0.85 处入页，向右上弯至 t_max 公式下方，箭头指向 t_max；页面下半部其余为空白横线（背面透字已忽略）
% 图为铅笔画的一族过同一点的圆，加红笔坐标轴、斜线、虚线、箭头和标注（均在图里）
%FIG p300_a 0.05 0.0 0.95 0.50
\notefig[0.85]{figures/p300_a.png}

\noindent\red{弦最长\unclear{为$\infty$}}\par
\noindent\red{直径当弦长}\par
\noindent\red{$\theta=2\alpha$\quad $t_{\max}=\dfrac{2\alpha m}{qB}$\quad 圆心角$=2\times$弦切角}
%%%END
%%%PAGEEND 300
